\subsection{Bounds on the norm of $\bar{B}^{\ell}$ and $R^{\ell,j}$.}
\label{sec:norm_bounds}

All of the ingredients are gathered to bound the matrix norms.

\begin{theorem}
\label{spec_B_bar}
Let 
$\ell = \lfloor c\log(n) \rfloor$ 
where $c < \frac{1}{32}$ 
is a universal constant. 
It holds that
\[
\|\bar{B}^{(\ell)}\|
\leq 
\left((d_1-1)(d_2-1)\right)^{\ell/4}
\exp( \log^{3/4} n )
\]
asymptotically almost surely.
\end{theorem}

\begin{proof}
The following holds for any natural number $k$, but for our proof, we will take
\begin{equation}
 k = \lfloor \log (n)^{1/3} \rfloor \quad \mbox{ and } \quad
 \ell = \lfloor c \log n \rfloor \mbox{ for some } c < \frac{1}{32}.
\end{equation}
We have
\begin{align}
\label{trace_expectation}
\E \left(\| \bar{B}^{(\ell)} \|^{2k} \right)
\leq 
\E
\left( 
\mathrm{Tr} 
 \left[ 
 \left(
 \bar{B}^{(\ell)} \bar{B}^{(\ell)^*} 
 \right)^{k} 
 \right]
\right)
=
\E\left( 
\sum_{\gamma \in \mathcal{C}} 
\prod_{t=1}^{2 k \ell} \bar{B}_{e_{2t-1} e_{2t+1}}
\right).
\end{align}
The sum is taken over the set $\mathcal{C}$ of all circuits 
$\gamma$
of length $2k\ell$,
where
$\gamma= ( \gamma_1, \gamma_2, \ldots, \gamma_{2k} )$ 
is formed by concatenation of $2k$ tangle-free segments
$\gamma_s \in F^{\ell}$, 
with the convention $e^{s+1}_1=e^s_{\ell+1}$.
Again, refer to Figure~\ref{fig:circuit} for clarification.

As in Section~\ref{sec:path_counting}, we will break these into circuits
which visit exactly 
$\nvertex = |V(\gamma)|$ different vertices, 
$\nvright = |V(\gamma) \cap V_2|$ of them in the right set, and
$\nedge = |E(\gamma)|$ different edges.
We define three disjoint sets of circuits:
\begin{align*}
%\begin{itemize}
	%\item 
    \mathcal{C}_1 & =
    \{\gamma \in \mathcal{C}:~\mbox{all edges in $\gamma$ are traversed at least twice} \}, \\
	%\item 
    \mathcal{C}_2 & =
    \{\gamma \in \mathcal{C}:~\mbox{at least one edge in $\gamma$ is traversed exactly once 
    and $\nvertex \leq kl+1$} \},~\mbox{and} \\
	%\item 
    \mathcal{C}_3 & =
    \{\gamma \in \mathcal{C}:~\mbox{at least one edge in $\gamma$ is traversed exactly once 
    and $\nvertex > kl+1$} \} .
\end{align*}
Define the quantities
\[
I_j
= 
\left|
\E
  \sum_{\gamma\in \mathcal{C}_j} 
  \prod_{t=1}^{2k\ell} \bar{B}_{e_{2t-1} e_{2t+1}} 
\right|
\leq
\sum_{\gamma \in \mathcal{C}_j}
\left|
\E
\prod_{t=1}^{2k\ell}
\bar{M}_{e_{2t-1} e_{2t}}
\right|
\]
for $j = 1, 2$ and 3, so that
\eqref{trace_expectation}
can be bounded as
\begin{align}\label{bound:I1+I2+I3}
\E\left(\|\bar{B}^{(\ell)}\|^{2k}\right)\leq I_1+I_2+I_3.
\end{align}

We will bound each term on the right hand side above.
The reason for this division is that, 
by Theorem~\ref{expectation_bound}, 
when we have any two-path traversed exactly once,
the expectation of the corresponding circuit is smaller,
because the matrix $\bar{B}$ is nearly centered.
We will see that the leading order terms in Eqn.~\eqref{trace_expectation} 
will come from circuits in $\mathcal{C}_1$. 
From Lemmas~\ref{expectation_bound} and
\ref{lem:circuit_count} and Corollary~\ref{cor:edge_iso}, we get that
\begin{align}
\nonumber
I_j
&\leq
\sum_{\nvertex,\nedge,\nvright}
\sum_{\gamma\in \mathcal{C}_j \cap \mathcal{C}_{\nvertex,\nedge}^\nvright} 
\left|
\mathbb{E} 
\prod_{t=1}^{2k\ell} \bar{M}_{e_{2t-1} e_{2t}}
\right|
\\
&\leq
\sum_{\nvertex,\nedge,\nvright}
C_{\nvertex,\nedge}^\nvright \,     
I_{\nvertex, \nedge}^\nvright \,
\left|
\mathbb{E} 
\prod_{t=1}^{2k\ell} \bar{M}_{e_{2t-1} e_{2t}}
\right|
\nonumber
\\
&\leq
\sum_{\nvertex,\nedge,\nvright}
n^{\nvertex - \nvright}
m^{\nvright}
(2\ell)^{4k} ((\nvertex + 1)^2 (\ell + 1))^{2k(\chi + 1)} \nonumber \\
& \qquad\qquad
\cdot
(d_1 (d_1 - 1))^{\nvertex - \nvright}
(d_2 (d_2 - 1))^\nvright
(d - 1)^{2(\chi - 1)} \nonumber \\
& \qquad\qquad
\cdot 
C
2^{b_\mathcal{C}} \left( \frac{1}{|E|} \right)^\nedge
\left( \frac{6 k \ell}{\sqrt{|E|}} \right)^{\nedge_1} \nonumber \\
&\leq
C
\sum_{\nvertex,\nedge,\nvright}
(2\ell)^{4k} ((\nvertex + 1)^2 (\ell + 1))^{2k(\chi + 1)}
(d_1 - 1)^{\nvertex - \nvright}
(d_2 - 1)^\nvright
(d - 1)^{2(\chi - 1)} \nonumber \\
& \qquad\qquad
\cdot 
2^{b_\mathcal{C}} 
\left( \frac{1}{|E|} \right)^{\nedge - \nvertex}
\left( \frac{6 k \ell}{\sqrt{|E|}} \right)^{\nedge_1} \nonumber \\
&\leq
C
\sum_{\nvertex,\nedge,\nvright}
(2\ell)^{4k} ((\nvertex + 1)^2 (\ell + 1))^{2k(\chi + 1)}
\left( (d_1 - 1)(d_2 - 1) \right)^{\nedge / 2}
\nonumber \\
& \qquad\qquad
\cdot
\left( \frac{d_2 - 1}{d_1 - 1} \right)^\psi
(d_1 - 1)^{1-\chi}
(d - 1)^{2(\chi - 1)}
2^{b_\mathcal{C}}
\left( \frac{1}{|E|} \right)^{\nedge - \nvertex}
\left( \frac{6 k \ell}{\sqrt{|E|}} \right)^{\nedge_1} 
\nonumber \\
&\leq
c_0 n
\ell^{4k} 
c_1^{k^2}
c_2^k
\sum_{\nvertex,\nedge}
\alpha^{2\nedge}
\nvertex
((\nvertex + 1)^2 (\ell + 1))^{2k(\chi + 1)}
\left( \frac{c_3}{n} \right)^\chi
\left( \frac{c_4 k \ell}{\sqrt{n}} \right)^{\nedge_1} 
\label{bound:I_j}
\end{align}
We use $C, c_0, c_1, c_2, c_3, c_4$ to denote constant terms
and set $\alpha = ((d_1-1)(d_2-1))^{1/4}$.
In the last line we used Lemma~\ref{lem:imbalance}
and Lemma~\ref{lem:inconsistent}
to bound 
$\psi$ and $b_\mathcal{C}$ 
in terms of 
$k$ and $\chi$
and remove the sum over $\nvright$,
which contains at most $\nvertex$ terms.
We will use Eqn.~\ref{bound:I_j} to bound each $I_j$. 

% Now we will separate the two sums and change variables to $\chi$,
% yielding \textcolor{red}{(check what bound we got for $b$.)}
% \begin{equation}
% I_j \leq
% c_0 n
% \ell^{4k} 
% c_1^{k^2}
% c_2^k
% \sum_{\nvertex}
% \alpha^{2 (\nvertex - 1)}
% \nvertex
% \sum_\chi
% ((\nvertex + 1)^2 (\ell + 1))^{2k(\chi + 1)}
% \left( \frac{c_3}{n} \right)^\chi
% \left( \frac{c_4 k \ell}{\sqrt{n}} \right)^{\nedge_1} .
%     \label{bound:I_j}
% \end{equation}
% \textcolor{red}{My only suggestion will be to omit equation 23.}
% \todo{I was thinking to absorb this into $c_3$}
% 

\subsubsection{Bounding $I_1$}

Here, $\nedge_1 = 0$ since every edge is traversed twice.
% In all cases, each circuit traverses $2k \ell$ \ps. 
% Hence, for each $\gamma\in \mathcal{C}_1$, 
% where each \2 is repeated twice,
% we have at most $k\ell$ different \ps. 
% Furthermore, since each edge can be in multiple \ps, 
% we have that the total number of different \ps 
% ~is greater than or equal to the total number of edges traversed by $\gamma$. 
We then have that 
$\nvertex - 1 \leq \nedge \leq k\ell$. 
Since $\gamma$ is connected,
we have 
$1 \leq \nvertex \leq k\ell+1$. 
Thus, on the right hand side of 
Eqn.~\eqref{bound:I_j} we get 
\begin{align*}
I_1 
&\leq 
c_0 n
\ell^{4k} 
c_1^{ k^2}
c_2^k
\sum_{\nvertex = 1}^{k\ell+1} 
\sum_{\nedge=\nvertex-1}^{k\ell} 
\alpha^{2\nedge}
\nvertex
((\nvertex + 1)^2 (\ell + 1))^{2k(\chi + 1)}
\left( \frac{c_3}{n} \right)^\chi
\\
&=
c_0 n
\ell^{4k} 
c_1^{ k^2}
c_2^k
\sum_{\nvertex = 1}^{k\ell+1} 
\alpha^{2(\nvertex - 1)}
\nvertex
((\nvertex + 1)^2 (\ell + 1))^{2k}
\sum_{\chi= 0}^{k\ell - \nvertex + 1} 
\left( 
    \alpha^2
    ((\nvertex + 1)^2 (\ell + 1))^{2k} 
    \frac{c_3}{n} 
\right)^\chi
\end{align*}
The second sum
is upper-bounded by 
$
\sum_{\chi=0}^\infty 
\left( C \frac{ ((\nvertex+1)^2 (\ell+1) )^{2k} }{n} \right)^\chi
$,
where $C$ is some constant,
which we show next is bounded by a common constant for our choices of $k$ and $\ell$ 
and all $\nvertex \in [k\ell + 1]$. 
To see this, it will suffice to show that 
$((\nvertex+1)^2 (\ell+1) )^{2k}=o(n)$. 
But $((\nvertex+1)^2 (\ell+1) ) = O(k^2 \ell^3)$
which, for our choices of $k$ and $\ell$, yields
\[
2k \log( (\nvertex+1)^2 (\ell+1) )=O({\log(n)^{1/3}}\log(\log(n)))=o(\log(n))
\]
as desired.

Finally, the first summand is maximized for $\nvertex = k\ell + 1$ and there are at $k\ell + 1$ many terms in that sum.
Therefore, modifying the constant $c_0$ yields
\begin{align}\label{bound:I_1}
I_1
\leq 
c_0' n 
\ell^{4k}
c_1^{ k^2}
c_2^k
(k \ell + 1)^2
\left( (k\ell+2)^2 (\ell+1) \right)^{2k}
\alpha^{2 k \ell}.
\end{align}

\subsubsection{Bounding $I_2$}

Here there is at least one edge traversed exactly once, so
we have $\nedge \geq \nvertex$ for $\gamma\in \mathcal{C}_2$. 
Taking $\nedge_1 = 0$ only increases the right hand side on
Eqn.~\eqref{bound:I_j}; it becomes
\begin{align*}
I_2
&\leq 
c_0 n
\ell^{4k} 
c_1^{ k^2}
c_2^k
\sum_{\nvertex=1}^{k\ell+1} \sum_{\nedge=\nvertex}^{2k\ell} 
\alpha^{2\nedge}
\nvertex
((\nvertex + 1)^2 (\ell + 1))^{2k(\chi + 1)}
\left( \frac{c_3}{n} \right)^\chi
\\
&=
c_0 n
\ell^{4k} 
c_1^{ k^2}
c_2^k
\sum_{\nvertex = 1}^{k\ell+1} 
\alpha^{2(\nvertex-1)}
\nvertex
((\nvertex + 1)^2 (\ell + 1))^{2k}
\sum_{\chi= 1}^{2k\ell - \nvertex + 1} 
\left( 
    \alpha^2
    ((\nvertex + 1)^2 (\ell + 1))^{2k} 
    \frac{c_3}{n} 
\right)^\chi
\end{align*}
Notice that this last term is almost identical to the one in the bound of $I_1$, 
except that now we start the second sum at $\chi=1$,
which leads to an extra factor of
$
O( ((\nvertex + 1)^2 (\ell + 1))^{2k} / n )
$.
This allow us to factor out another geometric series
and proceed as we did for $I_1$. 
This yields 
\begin{align}
I_2 
\leq 
c_0'
\ell^{4k}
c_1^{ k^2}
c_2^k
(k \ell+1)^2
\left( (k\ell+2)^2 (\ell+1) \right)^{4k}
\alpha^{2k\ell},
\label{bound:I_2}
\end{align} 
since there are $k\ell+1$ terms in the first sum.

\subsubsection{Bounding $I_3$}

This set will require more delicate treatment,
since circuits in $\mathcal{C}_3$ 
visit potentially many vertices and edges, 
yet we need to keep the power of $\alpha$ at most $2k\ell$.

We first show that, in this case, $\nedge_1$ is also large. 
We have 
$\nedge \geq \nvertex$,
and let $\nvertex = k\ell + t$.
Define
$\nedge_1'$ as the number of edges 
traversed once in $\gamma$, 
so that 
$\nedge_1' = b + \nedge_1$.
Since $\gamma$ has length $2k \ell$, we deduce that
$2 (\nedge - \nedge'_1) + \nedge'_1 \leq 2k\ell$,
which implies that $\nedge'_1 \geq 2t$.
Finally, Lemma~\ref{lem:inconsistent}
yields
$\nedge_1 \geq (2t - 4 (\chi + k))_+$.
Eqn.~\eqref{bound:I_j} then gives,
\begin{align*}
I_3 
&\leq 
c_0 n
\ell^{4k} 
c_1^{ k^2}
c_2^k
\sum_{\nvertex=k\ell+1}^{2k\ell} 
\sum_{\nedge=\nvertex}^{2k\ell} 
\alpha^{2\nedge}
\nvertex
((\nvertex + 1)^2 (\ell + 1))^{2k(\chi + 1)}
%& \qquad\qquad \cdot
\left( \frac{c_3}{n} \right)^\chi 
\left( \frac{(c_4 k \ell)^2}{n} \right)^
{(\nvertex - k\ell - 2(\chi+k))_+} 
\nonumber \\
&=
c_0 n
\ell^{4k} 
c_1^{ k^2}
c_2^k
\sum_{t=1}^{k\ell} 
\sum_{\chi=1}^{k\ell - t+1} 
\alpha^{2 (k\ell + \chi + t-1)}
(k\ell + t)
((k\ell + t + 1)^2 (\ell + 1))^{2k(\chi + 1)}\\ 
& \qquad\qquad\qquad\qquad\qquad\qquad\qquad\qquad\qquad\qquad \cdot
\left( \frac{c_3}{n} \right)^{\chi} 
\left( \frac{(c_4 k \ell)^2}{n} \right)^{(t - 2(\chi+k))_+}
\nonumber \\
&=
c_0
n
\ell^{4k} 
c_1^{ k^2}
c_2^k
\alpha^{2 (k\ell-1)}
% \nonumber \\
% & \qquad\qquad \cdot
\sum_{t=1}^{k\ell} 
(k\ell + t)
((k\ell + t + 1)^2 (\ell + 1))^{2k}
\alpha^{2t} 
\\
& \qquad\qquad\qquad\qquad\qquad\qquad \cdot
\sum_{\chi=1}^{k\ell - t+1}
\left( \alpha^2 ((k\ell + t + 1)^2 (\ell + 1))^{2k} \frac{c_3}{n} \right)^{\chi} 
% \left( C'
% \frac{ (C_{\nvertex}\ell)^{2k}}{n}
% \right)^{\chi} 
\left( \frac{(c_4 k \ell)^2}{n} \right)^{(t - 2(\chi+k))_+}.
\end{align*}

To simplify our notation, we will write 
\[
F(k,\ell)=c_0
n
\ell^{4k} 
c_1^{ k^2}
c_2^k
\alpha^{2 (k\ell-1)}
(2k\ell)
((2k\ell + 1)^2 (\ell + 1))^{2k}.
\]
Observe now that we can write:
\[
I_3\leq F(k,\ell)\sum_{t=1}^{k\ell} \alpha^{2t}\sum_{\chi=1}^{k\ell-t+1} \left(\frac{c(k,\ell)}{n}\right)^{\chi} \left(\frac{(c_4k\ell)^2}{n}\right)^{(t-2k-2\chi)_+}
\]
where $c(k,\ell)=c_3\alpha^2((2k\ell+1)^2(\ell+1))^{2k}$.
To bound the double sum on the right hand side above, 
we start by removing a factor of $\frac{c(k,\ell)}{n}$, 
whichs leaves
\begin{equation}\label{I_3bound}
I_3\leq F(k,\ell)\frac{c(k,\ell)}{n} \sum_{t=1}^{k\ell} \alpha^{2t}\sum_{\chi=0}^{k\ell-t} \left(\frac{c(k,\ell)}{n}\right)^{\chi} \left(\frac{(c_4k\ell)^2}{n}\right)^{(t-2k-2-2\chi)_+}
\end{equation}
The $n$ in the denominator is crucial to cancel the linear term in $F(k,\ell)$, keeping the upper bound for $I_3$ small. We focus on bounding the double sum. We split the sum in $t$ in two parts.\\
{\bf Case 1: $t< 2k+2$.} For these values of $t$, we have $(t-2k-2-2\chi)_+=0$, hence
\begin{equation}\label{I_3:eq1}
\sum_{t=1}^{2k+1} \alpha^{2t}\sum_{\chi=0}^{k\ell-t} \left(\frac{c(k,\ell)}{n}\right)^{\chi} \left(\frac{(c_4k\ell)^2}{n}\right)^{(t-2k-2-2\chi)_+}=\sum_{t=1}^{2k+1} \alpha^{2t}\sum_{\chi=0}^{k\ell-t} \left(\frac{c(k,\ell)}{n}\right)^{\chi}=O(\alpha^{4k})
\end{equation}
where the last equality uses again the same geometric upper bound for the sum over $\chi$.\\
{\bf Case 2: $t\geq 2k+2$.} 
We split the second sum, from $\chi=0$ to $N=\lfloor t/2-k-1\rfloor$ and
the terms with $\chi > N$ and analyse the two separately. 
The first can be upper bounded by
\[
\sum_{t=2k+2}^{k\ell} \alpha^{2t}\sum_{\chi=0}^{N} 
\left(\frac{c(k,\ell)}{n}\right)^{\chi} 
\left(\frac{c_4 k \ell}{\sqrt{n}}\right)^{(t/2-k-1-\chi)}
\leq
\alpha^{4k+4} \sum_{t=0}^{k\ell-2k-2}
\alpha^{2t} (N+1)
\left( \frac{c_4 k \ell}{\sqrt{n}} \right)^{\frac{t}{2}}.
\]
The last inequality can be checked in two steps: 
We first factor out the power of $\alpha^{4k+4}$, 
and then use that 
$\frac{c(k,\ell)}{n}\leq \frac{c_4 k \ell}{\sqrt{n}}$,
which holds for large enough $n$,
to simplify the second sum to the addition of $N+1$ equal terms. 
To bound the right hand side, 
we one more time upper bound by a 
geometric series of ratio less than one to get
\begin{equation}\label{I_3eq2}
\sum_{t=2k+2}^{k\ell} \alpha^{2t}\sum_{\chi=0}^{N} 
\left(\frac{c(k,\ell)}{n}\right)^{\chi} 
\left(\frac{c_4 k\ell}{\sqrt{n}}\right)^{(t/2-k-1-\chi)}
\leq
O(k\ell\alpha^{4k}).
\end{equation}
We are left with the terms $\chi > N$. In this case we get $(t-2k-2-2\chi)_+=0$, so
\[
\sum_{t=2k+2}^{k\ell} \alpha^{2t}\sum_{\chi=N+1}^{k\ell-t}
\left(\frac{c(k,\ell)}{n}\right)^{\chi} 
\left(\frac{c_4 k \ell}{\sqrt{n}}\right)^{(t/2-k-1-\chi)_+}
=
\sum_{t=2k+2}^{k\ell} \alpha^{2t}\sum_{\chi=N+1}^{k\ell-t} \left(\frac{c(k,\ell)}{n}\right)^{\chi}.
\]
The sum over $\chi$ is of the order of $\left(\frac{c(k,\ell)}{n}\right)^{N+1}\leq \left(\frac{c(k,\ell)}{n}\right)^{t/2-k-1}$. 
Substituting this into the above, we are left with
\[
\sum_{t=2k+2}^{k\ell} \alpha^{2t}\sum_{\chi=N+1}^{k\ell-t} \left(\frac{c(k,\ell)}{n}\right)^{\chi}
=
C\sum_{t=2k+2}^{k\ell} \alpha^{2t}\left(\frac{c(k,\ell)}{n}\right)^{t/2-k-1}
\]
for some universal constant $C$. After factoring $\alpha^{4k+4}$ and changing variables in the summation, we conclude that
\begin{equation}\label{I_3eq3}
\sum_{t=2k+2}^{k\ell} \alpha^{2t}\sum_{\chi=N+1}^{k\ell-t} \left(\frac{c(k,\ell)}{n}\right)^{\chi}
=
C\alpha^{4k+4}\sum_{t=0}^{k\ell-2k-2} \alpha^{2t}\left(\frac{c(k,\ell)}{n}\right)^{t/2}
=
O(\alpha^{4k}).
\end{equation}

Using \eqref{I_3bound} and the results for
case 1 \eqref{I_3:eq1} and case 2 (\ref{I_3eq2} and \ref{I_3eq3}),
we conclude that
\begin{equation}\label{I_3final_bound}
I_3\leq F(k,\ell)\frac{c(k,\ell)}{n}O(k\ell \alpha^{4k})
=
c_0'\ell^{4k} 
c_1^{ k^2}
(c_2')^k
(k\ell)^2
((2k\ell + 1)^2 (\ell + 1))^{4k}
\alpha^{2k\ell}
\end{equation}


\subsubsection{Finishing the proof of Theorem~\ref{spec_B_bar}}

We have bounded the three pieces we need to prove the theorem. 
From \eqref{bound:I_1}, \eqref{bound:I_2}, and \eqref{I_3final_bound},
with $n$ sufficiently large, we get
\begin{align*}
\E \left( \| \bar{B}^{(\ell)} \|^{2k}\right)
&\leq 
I_1 + I_2 + I_3  \\
&\leq 
\alpha^{2k\ell}
n
\cdot
C
c_1^{ k^2}
c_5^k
\ell^{4k}
(k \ell)^{4}
\left(k^2\ell^3 \right)^{4k}
\\
&\leq
\alpha^{2k\ell} 
n
\cdot
C'
c_1^{(\log n)^{2/3}}
c_6^{(\log n)^{1/3}}
(\log n)^{4 (\log n)^{1/3}}
(\log n)^{16/3}
\,
((\log n)^{11/3})^{4 (\log n)^{1/3}}
\\
&\leq
\alpha^{2k\ell} 
n
\cdot
C''
c_1^{(\log n)^{2/3}}
c_6^{(\log n)^{1/3}}
(\log n)^{20 (\log n)^{1/3} + 6} 
\\
& := 
\alpha^{2k\ell} n \cdot f(n) ,
\end{align*}
where $c_5, c_6, C, C', C''$ are universal constants.

Take any $\epsilon > 0$.
It can be checked that 
$(\log n)^{20 (\log n)^{1/3} + 6}= o(n^{\epsilon})$, and 
$f(n) = o(n^{\epsilon})$ as well.
Let $g(n) = \exp((\log n)^{3/4})$;
then $g(n) = o(n^\epsilon)$ but
$g(n)^{2k} \gg n^\epsilon$.
We apply Markov's inequality, so that
\begin{align}
\P \left[
\| B^{(\ell)} \| 
\geq 
\alpha^{\ell} g(n)
\right]
&\leq 
\frac{ \E \left( \| \bar{B}^{(\ell)} \|^{2k} \right) }{ \alpha^{2k\ell} g(n)^{2k}}
\leq  
\frac{n f(n)}{g(n)^{2k}}=o(1),
\label{eq:f_over_g}
\end{align}
which is the statement of the theorem.
\end{proof}

\begin{theorem}
\label{spec_R}
Let $1\leq j\leq \ell = \lfloor c \log(n) \rfloor$ 
where $c < \frac{1}{32}$ is a universal constant.
Then
\[
\|R^{\ell,j}\| 
\leq 
\frac{(d-1)^\ell}{n} \exp((\log n)^{3/4})
\]
asymptotically almost surely.
\end{theorem}

\begin{proof}
The proof is analogous to the proof of Theorem \ref{spec_B_bar}. Recall the definition of $R^{\ell,j}$ in Eqn.~\ref{error_term}.
For any integer $k$, we have that
\begin{align}
\E\left(\|R^{\ell,j}\|^{2k}\right)
&\leq 
\E 
\left(
\mathrm{Tr}
\left[
 \left( (R^{\ell,j}) (R^{\ell,j})^* \right)^{k}
\right]
\right)
\nonumber \\
&\leq
\sum_{ \gamma \in \mathcal{D} }
\prod_{s=1}^{2k}
\left|
\E
\prod_{i=1}^{j-1} \bar{B}_{e^s_{2i-1} e^s_{2i+1}} 
S_{e^s_{2j-1} e^s_{2j+1}}
\prod_{i=j+1}^{\ell} B_{e^s_{2i-1} e^s_{2i+1}} 
\right| 
\nonumber \\
&=
\sum_{ \gamma \in \mathcal{D} }
\prod_{s=1}^{2k}
\left|
\E
\prod_{i=1}^{j-1} \bar{M}_{e^s_{2i-1} e^s_{2i}} 
S_{e^s_{2j-1} e^s_{2j+1}}
\prod_{i=j+1}^{\ell} M_{e^s_{2i-1} e^s_{2i}} 
\right| 
\label{trace_expectation:R^l,j}
\end{align}
Now, the sum is over 
the set $\mathcal{D}$ of circuits 
$\gamma= (\gamma_1,\gamma_2,\dots,\gamma_{2k})$ of length $2k\ell$
formed from $2k$ elements of $T^{\ell,j}$, 
$\gamma_s = (e^s_1, e^s_2, \ldots, e^s_{2\ell+1} )$ for $s \in [2k]$,
with the convention $e^{s+1}_1 = e^s_{2\ell+1}$ (see definition ~\ref{def:tangled_paths}).


% We proceed as before.
% Let
% \begin{equation}
%     K_j = \sum_{\gamma \in \mathcal{D}_j}
%     \prod_{s=1}^{2k}
%     \left|
%     \E
%     \prod_{i=1}^{j-1} \bar{M}_{e^s_{2i-1} e^s_{2i}} 
%     S_{e^s_{2j-1} e^s_{2j+1}}
%     \prod_{i=j+1}^{\ell} M_{e^s_{2i-1} e^s_{2i}} 
%     \right| .
% \end{equation}
% Then,
% \begin{equation}
%     \E\left(\|R^{\ell,j}\|^{2k}\right)
%     \leq
%     K_1 + K_2 + K_3.
% \end{equation}
% \todo{Hmm... above probably not necessary}

Proceeding as before, using Lemma \ref{lem:combinatorial_of_reminder} and Corollary \ref{cor:edge_iso} we get that
\begin{align*}
    \E\left(\|R^{\ell,j}\|^{2k}\right)
    &\leq
    \sum_{\nvertex, \nedge, \nvright}
    \sum_{\gamma \in \mathcal{D}_{\nvertex, \nedge}^\nvright }
    \prod_{s=1}^{2k}
    \left|
    \E
    \prod_{i=1}^{j-1} \bar{M}_{e^s_{2i-1} e^s_{2i}} 
    S_{e^s_{2j-1} e^s_{2j+1}}
    \prod_{i=j+1}^{\ell} M_{e^s_{2i-1} e^s_{2i}}
    \right|  \\
    &\leq
    \sum_{\nvertex, \nedge, \nvright}
    D_{\nvertex, \nedge}^\nvright 
    I_{\nvertex, \nedge}^\nvright 
    \prod_{s=1}^{2k}
    \left|
    \E
    \prod_{i=1}^{j-1} \bar{M}_{e^s_{2i-1} e^s_{2i}} 
    S_{e^s_{2j-1} e^s_{2j+1}}
    \prod_{i=j+1}^{\ell} M_{e^s_{2i-1} e^s_{2i}}
    \right| \\
    &\leq
    \sum_{\nvertex, \nedge, \nvright}
    m^\nvright n^{\nvertex-\nvright}
    (2 \ell)^{6k} 
    ( (\nvertex + 1)^2 (\ell + 1) )^{6k(\chi + 1)} \\
    & \qquad \qquad \cdot
    (d_1 (d_1 - 1))^{\nvertex - \nvright}
    (d_2 (d_2 - 1))^\nvright
    (d - 1)^{2 (\chi - 1)} \\
    & \qquad \qquad \cdot
    \prod_{s=1}^{2k}
    \left( \frac{d-1}{|E|} \right)
    C
    2^{b_s}
    \left(\frac{1}{|E|} \right)^{\nedge_s}
    \left( \frac{6\ell}{\sqrt{|E|}} \right)^{\nedge_{1,s}} .
\end{align*}
The last inequality uses Lemma \ref{expectation_bound} on each path $\gamma_s$, 
with  $\nedge_s$, $b_s$, and $\nedge_{1,s}$ 
defined analogous to the quantities $\nedge$, $b$ and $\nedge_1$ in the same lemma.
Note that $S_{ef} \leq \frac{d-1}{|E|}$ for any $e,f$ and 
$\sum_{s=1}^{2k} b_s \leq b_\mathcal{D} \leq 16 (k + \chi)$
by Lemma~\ref{lem:inconsistent_tangled}. 
Setting
$\nedge = \sum_{s=1}^{2k} \nedge_s$, 
taking
$\nedge_1 = \sum_{s=1}^{2k} \nedge_{1,s} \geq 0$,
using $d_1, d_2 \leq d$,
and combining terms, 
we find that
\begin{align*}
    \E\left(\|R^{\ell,j}\|^{2k}\right)
    & \leq
    \sum_{\nvertex, \nedge, \nvright}
    (2 \ell)^{6k} 
    ( (\nvertex + 1)^2 (\ell + 1) )^{6k(\chi + 1)} \\
    & \qquad \qquad \cdot
    (d_1 - 1)^{\nvertex - \nvright}
    (d_2 - 1)^\nvright
    (d - 1)^{2 (\chi - 1)} \\
    & \qquad \qquad \cdot
    \left( \frac{C'}{|E|} \right)^{2k}
    2^{b_\mathcal{D}}
    \left(\frac{1}{|E|} \right)^{\nedge - \nvertex} \\
    &\leq
    c_0 
    \ell^{6k} 
    c_1^k 
    n^{1 - 2k}
    \sum_{\nvertex, \nedge}
    \nvertex
    ( (\nvertex + 1)^2 (\ell + 1) )^{6k(\chi + 1)}
    (d - 1)^{\nvertex}
    \left(
    \frac{c_2}{n}
    \right)^{\chi} \\
    &\leq
    c_0 
    \ell^{6k} 
    c_1^k 
    n^{1 - 2k}
    \sum_{\nvertex = 1}^{2k\ell}
    \sum_{\nedge = \nvertex - 1}^{2k\ell}
    \nvertex
    ( (\nvertex + 1)^2 (\ell + 1) )^{6k(\chi + 1)}
    (d - 1)^{\nvertex}
    \left(
    \frac{c_2}{n}
    \right)^{\chi} \\
    &\leq
    c_0 
    \ell^{6k} 
    c_1^k 
    n^{1 - 2k}
    \sum_{\nvertex = 1}^{2k\ell}
    \nvertex
    ( (\nvertex + 1)^2 (\ell + 1) )^{6k}
    (d - 1)^{\nvertex}
    \sum_{\chi = 0}^{2k\ell}
    \left(
    ( (\nvertex + 1)^2 (\ell + 1) )^{6k}
    \frac{c_2}{n}
    \right)^{\chi}
    ,
    % & \leq
    % (2 k \ell)^3
    % c_0 
    % \ell^{6k} 
    % c_1^k 
    % n^{1-2k}
    % ( (2 k \ell + 1)^2 (\ell + 1) )^{6k}
    % (d - 1)^{2 k \ell} 
    % \\
    % & = 
    % (k \ell)^3
    % c_3 
    % \ell^{6k} 
    % c_1^k 
    % n
    % ( (2 k \ell + 1)^2 (\ell + 1) )^{12 k^2 \ell}
    % \left(
    %     \frac{(d - 1)^\ell}{n}
    % \right)^{2 k}
\end{align*}
for some constants $c_0, c_1, c_2$.
Note that the important $n^{-2k}$ comes from
the $|E|^{-2k}$ deterministic term.
Now, like before we have that $( (\nvertex + 1)^2 (\ell + 1) )^{6k} = o(n)$,
so for sufficiently large $n$ we can bound the sum over $\chi$
by a constant.
The sum over $\nvertex$ is easily bounded as before, leading to
\begin{align*}
    \E\left(\|R^{\ell,j}\|^{2k}\right)
    &\leq 
    (d-1)^{2k\ell}
    n^{1-2k}
    \cdot
    c_0'
    \ell^{6k} 
    c_1^k 
    (2k\ell)^2
    ((2k\ell + 1)^2 (\ell + 1))^{6k}
    \\
    &\leq
    (d-1)^{2k\ell}
    n^{1-2k}
    \cdot
    c_0''
    c_1'^{(\log n)^{1/3}}
    (\log n)^{6(\log n)^{1/3}}
    (\log n)^{8/3}
    ((\log n)^{11/3})^{6(\log n)^{1/3}} \\
    &:= 
    (d-1)^{2k\ell}
    n^{1-2k}
    \cdot 
    f(n).
\end{align*}
To finish the proof, we note that $f(n) = o(n^\epsilon)$ 
for any $\epsilon > 0$.
Take $g(n) = \exp((\log n)^{3/4})$.
Then, applying Markov's inequality,
\begin{equation}
    \P 
    \left[ 
    \| R^{\ell,j} \| 
    \geq \frac{(d - 1)^{\ell} }{n} g(n)
    \right]
    \leq
    \frac{\E \left( \| R^{\ell,j} \|^{2k} \right) n^{2k}}
    { (d-1)^{2k\ell} g(n)^{2k} }
    \leq
    \frac{n f(n)}{g(n)^{2k}} = o(1).
\end{equation}

\subsection{Proof of the main result, Theorem~\ref{thm:gap_B}}
We will again take $\ell = \lfloor c \log n \rfloor$, with $c$ chosen so that the
graph is $\ell$-tangle-free 
with high probability.
By Eqns.~\eqref{eq:upper_bound_lambda2} and \eqref{decomposition},
\[
|\lambda_2|^{\ell}\leq \| \bar{B}^{(\ell)}\|+
\sum_{k=1}^{\ell} \| R^{\ell,j}\|.
\]
Notice that
$
(d-1)^\ell \leq (d-1)^{c \log n} \leq  n^{c \log d}
$,
so take 
$c < \min \left(\frac{1}{32}, \frac{1}{\log d} \right)$.
Then $(d-1)^\ell = O(n^\epsilon)$ for some $0 < \epsilon < 1$,
and recall that 
$\exp((\log n)^{3/4}) = o(n^{\epsilon'})$ for any 
$\epsilon' > 0$.
We apply
Theorems \ref{spec_B_bar} and \ref{spec_R} to get
\begin{eqnarray}
\nonumber
|\lambda_2| 
& \leq & 
\left(
\exp((\log n)^{3/4}) \left( (d_1-1)(d_2-1) \right)^{\ell/4} 
+ 
\ell 
\exp((\log n)^{3/4})
\frac{(d-1)^\ell}{n}
\right)^{1/\ell} \\
\nonumber
& = & \left( (d_1-1)(d_2-1) \right)^{1/4}+\epsilon_n,
\end{eqnarray}
where $\epsilon_n \to 0$ as $n \to \infty$.
\end{proof}